% Generated draft: requires mathematical and source review.
\documentclass[11pt]{article}
\usepackage[margin=25mm]{geometry}
\usepackage{fontspec}
\setmainfont{Latin Modern Roman}
\usepackage{amsmath,amssymb,mathtools,unicode-math}
\setmathfont{Latin Modern Math}
\usepackage{enumitem,needspace,xurl,hyperref}
\setlength{\parindent}{0pt}
\setlength{\parskip}{5pt}
\setlength{\emergencystretch}{3em}
\newcommand{\N}{\mathbb{N}}
\newcommand{\Z}{\mathbb{Z}}
\newcommand{\Q}{\mathbb{Q}}
\newcommand{\R}{\mathbb{R}}
\newcommand{\C}{\mathbb{C}}
\newcommand{\F}{\mathbb{F}}
\newcommand{\A}{\mathbb{A}}
\newcommand{\PP}{\mathbb{P}}
\newcommand{\E}{\mathbb{E}}
\newcommand{\eps}{\varepsilon}
\newcommand{\dd}{\,\mathrm{d}}
\begin{document}
\section*{UnsolvedMath ID: 30004640}
\textbf{Source problem code:} catalogue entry
\subsection{Definitions and mathematical statement}
A hyperbolic balance law is an evolution equation or system involving a hyperbolic flux and a source term. A fully discrete approximation discretizes both space and time. Let \(U_{h,k}\) denote a numerical approximation with spatial mesh parameter \(h>0\) and time step \(k>0\), and let \(u_{h,k}\) denote a corresponding reconstruction as a function of space and time. The reconstruction, the underlying balance law, and its data are not specified in the supplied record.

\textbf{Requested convergence problem.} Prove convergence of fully discrete approximations of hyperbolic balance laws, including the Lax--Friedrichs and Godunov schemes. Schematically, the intended conclusion is
\[
 u_{h,k}\longrightarrow u\qquad\text{as }h,k\to0,
\]
where \(u\) is an appropriate solution of the prescribed balance law, and the meshes and time steps satisfy the conditions required by the intended numerical scheme.

\textbf{Unresolved formulation.} This display is a statement template, not a fully quantified conjecture. The record does not specify the spatial dimension or domain, the number of unknowns, the flux and source, the precise hyperbolicity assumptions, the class or size of initial data, boundary conditions, or the time interval. It also does not specify an entropy or other admissibility condition for the limiting solution, the treatment of the source term, the exact fully discrete updates, mesh restrictions or a Courant--Friedrichs--Lewy condition, the reconstruction, or the topology of convergence. It is unclear whether convergence of the entire approximation family or only of subsequences is intended, and whether the target is a unique admissible solution. Consequently, the supplied record does not determine a precise universal convergence assertion without additional hypotheses.

The supplied dated literature assessment labels the problem partially solved: it reports convergence for selected \(2\times2\) systems, semidiscrete schemes, and classes with total-variation bounds, while describing general fully discrete Lax--Friedrichs or Godunov convergence as open. These are supplied qualifications, not independently verified results or a specification of the intended hypotheses.
\subsection{Short English statement}
The problem asks for convergence of fully discrete numerical approximations of hyperbolic balance laws, including Lax--Friedrichs and Godunov schemes. The supplied record leaves the equations, data, numerical assumptions, limiting solution concept, and convergence topology unspecified, so it supports only a qualified statement template. Its partially solved status and descriptions of known cases are supplied dated curation requiring independent review.
\subsection{Sources}
\begin{itemize}
\item Supplied local UnsolvedMath registry record, permanent ID \(30004640\), title \emph{Convergence of Fully Discrete Hyperbolic Balance-Law Schemes}, source problem code \texttt{catalogue entry}, in the Oberwolfach Reports Open Problems collection. The registry labels the problem \texttt{partially\_solved} and gives proposed year 2021; no individual proposer is supplied.
\item Supplied report citation: \emph{Hyperbolic Balance Laws: modeling, analysis, and numerics} (2022), Oberwolfach Reports, DOI 10.4171/owr/2021/11, pp.\ 589--661. \url{https://doi.org/10.4171/owr/2021/11}. The background describes Oberwolfach Reports as documenting mathematical workshops and their open-problem discussions. No precise page or passage identifying this problem is supplied.
\item Supplied literature source: Bressan--Chiri--Shen. \url{https://par.nsf.gov/servlets/purl/10249103}. No title, publication details, or theorem identifying the applicable convergence result is supplied.
\item Supplied literature source: Oberwolfach Report 11/2021. \url{https://doi.org/10.4171/owr/2021/11}.
\item Literature assessment checked at \texttt{2026-08-22}: this is supplied dated curation derived from the supplied OWR curation, not an independently verified assessment. It reports convergence for selected \(2\times2\) systems, semidiscrete schemes, and classes with total-variation bounds, while describing general fully discrete Lax--Friedrichs or Godunov convergence as open. Time-sensitive status and cited claims require independent verification.
\item The supplied statement status is \texttt{corrected\_verified}. Its accompanying verification describes a publication-format audit that preserved mathematical content while repairing notation or typography; it does not establish mathematical correctness or resolve the missing hypotheses.
\end{itemize}
\end{document}
